% 第十章：行波法

\chapter{行波法}

\begin{wulibj}[本章导读]
行波法是求解波动方程初值问题的直接方法。它利用波动方程的特征线，将偏微分方程转化为常微分方程，从而得到解的显式表达式——达朗贝尔公式。这种方法不仅给出了解的具体形式，还清晰地揭示了波的传播特性。
\end{wulibj}

% ========== 第一节 ==========
\section{达朗贝尔公式}

\subsection{一维波动方程的初值问题}

考虑无限长弦的自由振动：
\begin{equation}
    \begin{cases}
        \dfrac{\partial^2 u}{\partial t^2} = a^2\dfrac{\partial^2 u}{\partial x^2}, & x \in \mathbb{R}, \; t > 0 \\
        u(x, 0) = \varphi(x), & x \in \mathbb{R} \\
        \dfrac{\partial u}{\partial t}(x, 0) = \psi(x), & x \in \mathbb{R}
    \end{cases}
\end{equation}

\subsection{特征坐标变换}

关键思想：引入特征坐标，将方程化简。

特征线方程：$\frac{\mathrm{d}x}{\mathrm{d}t} = \pm a$，即 $x \mp at = C$。

令
\begin{equation}
    \xi = x - at, \quad \eta = x + at
\end{equation}

则
\begin{align}
    \frac{\partial u}{\partial x} &= \frac{\partial u}{\partial\xi}\frac{\partial\xi}{\partial x} + \frac{\partial u}{\partial\eta}\frac{\partial\eta}{\partial x} = \frac{\partial u}{\partial\xi} + \frac{\partial u}{\partial\eta} \\
    \frac{\partial^2 u}{\partial x^2} &= \frac{\partial^2 u}{\partial\xi^2} + 2\frac{\partial^2 u}{\partial\xi\partial\eta} + \frac{\partial^2 u}{\partial\eta^2}
\end{align}

类似地
\begin{equation}
    \frac{\partial^2 u}{\partial t^2} = a^2\left(\frac{\partial^2 u}{\partial\xi^2} - 2\frac{\partial^2 u}{\partial\xi\partial\eta} + \frac{\partial^2 u}{\partial\eta^2}\right)
\end{equation}

代入波动方程：
\begin{equation}
    a^2\left(\frac{\partial^2 u}{\partial\xi^2} - 2\frac{\partial^2 u}{\partial\xi\partial\eta} + \frac{\partial^2 u}{\partial\eta^2}\right) = a^2\left(\frac{\partial^2 u}{\partial\xi^2} + 2\frac{\partial^2 u}{\partial\xi\partial\eta} + \frac{\partial^2 u}{\partial\eta^2}\right)
\end{equation}

化简得
\begin{equation}
    \frac{\partial^2 u}{\partial\xi\partial\eta} = 0
\end{equation}

\subsection{通解与达朗贝尔公式}

\begin{dingli}{波动方程的通解}{}
一维波动方程 $\frac{\partial^2 u}{\partial t^2} = a^2\frac{\partial^2 u}{\partial x^2}$ 的通解为
\begin{equation}
    u(x, t) = f(x - at) + g(x + at)
\end{equation}
其中 $f$ 和 $g$ 是任意二次可微函数。

\begin{proof}
由 $\frac{\partial^2 u}{\partial\xi\partial\eta} = 0$，先对 $\xi$ 积分：
\begin{equation}
    \frac{\partial u}{\partial\eta} = h(\eta)
\end{equation}
再对 $\eta$ 积分：
\begin{equation}
    u = \int h(\eta)\,\mathrm{d}\eta + f(\xi) = g(\eta) + f(\xi) = g(x + at) + f(x - at)
\end{equation}
\end{proof}
\end{dingli}

\begin{tishi}[物理意义]
\begin{itemize}
    \item $f(x - at)$：向右传播的波（右行波），波形 $f$ 以速度 $a$ 向右移动
    \item $g(x + at)$：向左传播的波（左行波），波形 $g$ 以速度 $a$ 向左移动
\end{itemize}
波动方程的解是左右行波的叠加！
\end{tishi}

\begin{figure}[htbp]
\centering
\begin{tikzpicture}[scale=0.8]
    % 右行波
    \begin{scope}
        \draw[->, gray] (-0.5, 0) -- (5, 0) node[right] {$x$};
        \draw[->, gray] (0, -0.5) -- (0, 2) node[above] {$u$};
        
        \draw[blue, thick, domain=0:4.5, samples=50] plot (\x, {0.8*exp(-(\x-1)^2)});
        \draw[red, thick, dashed, domain=0:4.5, samples=50] plot (\x, {0.8*exp(-(\x-3)^2)});
        
        \draw[->, thick] (2, 1.5) -- (3, 1.5) node[midway, above] {$a$};
        
        \node[blue] at (1, -0.5) {$t=0$};
        \node[red] at (3, -0.5) {$t=t_1$};
        \node at (2.5, 2.3) {右行波 $f(x-at)$};
    \end{scope}
    
    % 左行波
    \begin{scope}[xshift=7cm]
        \draw[->, gray] (-0.5, 0) -- (5, 0) node[right] {$x$};
        \draw[->, gray] (0, -0.5) -- (0, 2) node[above] {$u$};
        
        \draw[blue, thick, domain=0:4.5, samples=50] plot (\x, {0.8*exp(-(\x-3.5)^2)});
        \draw[red, thick, dashed, domain=0:4.5, samples=50] plot (\x, {0.8*exp(-(\x-1.5)^2)});
        
        \draw[->, thick] (3, 1.5) -- (2, 1.5) node[midway, above] {$a$};
        
        \node[blue] at (3.5, -0.5) {$t=0$};
        \node[red] at (1.5, -0.5) {$t=t_1$};
        \node at (2.5, 2.3) {左行波 $g(x+at)$};
    \end{scope}
\end{tikzpicture}
\caption{行波传播示意图}
\label{fig:traveling-waves}
\end{figure}

\begin{dingli}{达朗贝尔公式}{}
利用初始条件确定 $f$ 和 $g$：
\begin{align}
    u(x, 0) &= f(x) + g(x) = \varphi(x) \\
    \frac{\partial u}{\partial t}(x, 0) &= -af'(x) + ag'(x) = \psi(x)
\end{align}

从第一式：$f(x) + g(x) = \varphi(x)$。

对第二式积分：
\begin{equation}
    -af(x) + ag(x) = \int_0^x \psi(s)\,\mathrm{d}s + C
\end{equation}

解得
\begin{align}
    f(x) &= \frac{1}{2}\varphi(x) - \frac{1}{2a}\int_0^x \psi(s)\,\mathrm{d}s - \frac{C}{2a} \\
    g(x) &= \frac{1}{2}\varphi(x) + \frac{1}{2a}\int_0^x \psi(s)\,\mathrm{d}s + \frac{C}{2a}
\end{align}

代入通解，常数 $C$ 被消去，得一维波动方程初值问题的解：
\[
    u(x, t) = \frac{1}{2}[\varphi(x-at) + \varphi(x+at)] + \frac{1}{2a}\int_{x-at}^{x+at} \psi(s)\,\mathrm{d}s
\]
\end{dingli}

\begin{wulibj}[达朗贝尔公式的物理意义]
初始位移 $\varphi(x)$ 分裂成两半，分别向左右传播。初始速度 $\psi(x)$ 的影响区域是 $[x-at, x+at]$，表示速度对点 $(x, t)$ 的影响来自特征线内的所有点。
\end{wulibj}

\begin{lizi}{}{}
求解无限长弦的振动问题：
\begin{equation}
    \begin{cases}
        u_{tt} = a^2 u_{xx}, & x \in \mathbb{R}, \; t > 0 \\
        u(x, 0) = \mathrm{e}^{-x^2}, & x \in \mathbb{R} \\
        u_t(x, 0) = 0, & x \in \mathbb{R}
    \end{cases}
\end{equation}

\begin{jie}
    这是典型的达朗贝尔公式应用。这里 $\varphi(x) = \mathrm{e}^{-x^2}$，$\psi(x) = 0$。

    代入达朗贝尔公式：
    \begin{equation}
        u(x, t) = \frac{1}{2}[\varphi(x-at) + \varphi(x+at)] + \frac{1}{2a}\int_{x-at}^{x+at} \psi(s)\,\mathrm{d}s
    \end{equation}

    由于 $\psi(x) = 0$，积分项为零：
    \begin{equation}
        u(x, t) = \frac{1}{2}\left[\mathrm{e}^{-(x-at)^2} + \mathrm{e}^{-(x+at)^2}\right]
    \end{equation}

    \textbf{物理意义}：初始时刻的高斯脉冲分裂成两个相同的脉冲，分别以速度 $a$ 向左右传播。每个脉冲的振幅是原来的一半。
\end{jie}
\end{lizi}

% ========== 第二节 ==========
\section{依赖区间与影响区域}

\subsection{依赖区间}

\begin{dingliyi}{依赖区间}{}
点 $(x_0, t_0)$ 的\textbf{依赖区间}是 $x$ 轴上的区间 $[x_0 - at_0, x_0 + at_0]$，即过该点的两条特征线与 $x$ 轴的交点之间的区域。
\end{dingliyi}

由达朗贝尔公式，$u(x_0, t_0)$ 完全由 $\varphi$ 和 $\psi$ 在依赖区间上的值决定。

\begin{figure}[htbp]
\centering
\begin{tikzpicture}[scale=1.0]
    % 坐标轴
    \draw[->, gray] (-0.5, 0) -- (5, 0) node[right] {$x$};
    \draw[->, gray] (0, -0.3) -- (0, 3.5) node[above] {$t$};
    
    % 特征线
    \draw[blue, thick] (1, 0) -- (3, 2);
    \draw[blue, thick] (4, 0) -- (2, 2);
    
    % 点
    \fill[red] (3, 2) circle (2pt) node[above right] {$(x_0, t_0)$};
    
    % 依赖区间
    \draw[green!60!black, thick] (1, 0) -- (4, 0);
    \fill[green!60!black, opacity=0.3] (1, 0) -- (3, 2) -- (4, 0) -- cycle;
    
    \fill (1, 0) circle (2pt) node[below] {$x_0-at_0$};
    \fill (4, 0) circle (2pt) node[below] {$x_0+at_0$};
    
    \node[green!60!black] at (2.5, 0.8) {依赖区间};
\end{tikzpicture}
\caption{依赖区间}
\label{fig:domain-of-dependence}
\end{figure}

\subsection{影响区域}

\begin{dingliyi}{影响区域}{}
初始时刻区间 $[x_1, x_2]$ 上的初始条件所影响的 $(x, t)$ 区域，称为\textbf{影响区域}。
\end{dingliyi}

初始时刻区间 $[x_1, x_2]$ 的影响区域是由过 $x_1$ 的右行特征线和过 $x_2$ 的左行特征线围成的三角形区域。

\begin{figure}[htbp]
\centering
\begin{tikzpicture}[scale=1.0]
    % 坐标轴
    \draw[->, gray] (-0.5, 0) -- (6, 0) node[right] {$x$};
    \draw[->, gray] (0, -0.3) -- (0, 3.5) node[above] {$t$};
    
    % 初始区间
    \draw[green!60!black, thick] (1, 0) -- (3, 0);
    \fill[green!60!black, opacity=0.3] (1, 0) -- (1, 0.1) -- (3, 0.1) -- (3, 0) -- cycle;
    
    % 特征线
    \draw[blue, thick] (1, 0) -- (4, 3);
    \draw[blue, thick] (3, 0) -- (0, 3);
    
    % 影响区域
    \fill[red, opacity=0.2] (1, 0) -- (4, 3) -- (0, 3) -- (3, 0) -- cycle;
    
    \fill (1, 0) circle (2pt) node[below] {$x_1$};
    \fill (3, 0) circle (2pt) node[below] {$x_2$};
    
    \node[red] at (2, 1.5) {影响区域};
\end{tikzpicture}
\caption{影响区域}
\label{fig:range-of-influence}
\end{figure}

\subsection{波的有限传播速度}

达朗贝尔公式揭示了波的一个重要性质：\textbf{有限传播速度}。

\begin{dingli}{惠更斯原理（波的有限传播速度）}{}
初始扰动只能在有限时间内传播有限距离。在时刻 $t$，初始时刻位于 $x_0$ 的扰动只能影响到 $|x - x_0| \leq at$ 的区域。
\end{dingli}

\begin{proof}
由达朗贝尔公式
\begin{equation}
    u(x, t) = \frac{1}{2}[\varphi(x-at) + \varphi(x+at)] + \frac{1}{2a}\int_{x-at}^{x+at} \psi(\xi)\,\mathrm{d}\xi
\end{equation}
解在点 $(x, t)$ 的值只依赖于初始数据在区间 $[x-at, x+at]$ 上的值。

因此，初始时刻位于 $x_0$ 的扰动（即 $\varphi(x_0)$ 和 $\psi(x_0)$），只有在 $x_0 \in [x-at, x+at]$ 时才能影响 $u(x, t)$。等价地，
\begin{equation}
    x - at \leq x_0 \leq x + at \quad \Longleftrightarrow \quad |x - x_0| \leq at
\end{equation}
这就是惠更斯原理的数学表述。
\end{proof}

\begin{lizi}{}{}
设初始位移 $\varphi(x)$ 仅在区间 $[2, 4]$ 内非零，初始速度 $\psi(x) = 0$。求 $t = 1$ 时波动到达的区域（设波速 $a = 1$）。

\begin{jie}
    由影响区域的定义，初始扰动在 $[x_1, x_2] = [2, 4]$ 内的影响区域是由过 $x_1 = 2$ 的右行特征线和过 $x_2 = 4$ 的左行特征线围成的区域。

    特征线方程：$x - at = C$（右行），$x + at = C'$（左行）。

    过 $x_1 = 2$ 的右行特征线：$x - t = 2$，即 $x = t + 2$。
    
    过 $x_2 = 4$ 的左行特征线：$x + t = 4$，即 $x = 4 - t$。

    当 $t = 1$ 时：
    \begin{itemize}
        \item 右行特征线到达 $x = 3$
        \item 左行特征线到达 $x = 3$
    \end{itemize}
    
    因此，$t = 1$ 时波动到达的区域是 $x \in [1, 5]$。

    \textbf{验证}：由达朗贝尔公式，$u(x, 1) = \frac{1}{2}[\varphi(x-1) + \varphi(x+1)]$。要使 $u(x, 1) \neq 0$，需要 $x-1 \in [2, 4]$ 或 $x+1 \in [2, 4]$，即 $x \in [1, 5]$。
\end{jie}
\end{lizi}

% ========== 第三节 ==========
\section{半无界弦的振动}

\subsection{延拓法}

对于半无界问题 $x > 0$，边界条件 $u(0, t) = 0$（固定端），可以用\textbf{奇延拓}方法求解。

\textbf{思路}：将半无界问题延拓为全无界问题，使边界条件自动满足。

对于 $u(0, t) = 0$，将初始条件作奇延拓：
\begin{equation}
    \tilde{\varphi}(x) = \begin{cases}
        \varphi(x), & x \geq 0 \\
        -\varphi(-x), & x < 0
    \end{cases}, \quad
    \tilde{\psi}(x) = \begin{cases}
        \psi(x), & x \geq 0 \\
        -\psi(-x), & x < 0
    \end{cases}
\end{equation}

然后用达朗贝尔公式求解。

\begin{lizi}{}{}
求解半无界弦振动问题：
\begin{equation}
    \begin{cases}
        \dfrac{\partial^2 u}{\partial t^2} = a^2\dfrac{\partial^2 u}{\partial x^2}, & x > 0, \; t > 0 \\
        u(x, 0) = \varphi(x), & x > 0 \\
        \dfrac{\partial u}{\partial t}(x, 0) = 0, & x > 0 \\
        u(0, t) = 0, & t > 0
    \end{cases}
\end{equation}

\begin{jie}
    将初始条件作奇延拓：$\tilde{\varphi}(x) = \varphi(x)$（$x \geq 0$），$\tilde{\varphi}(x) = -\varphi(-x)$（$x < 0$）。
    
    由达朗贝尔公式：
    \begin{equation}
        u(x, t) = \frac{1}{2}[\tilde{\varphi}(x-at) + \tilde{\varphi}(x+at)]
    \end{equation}
    
    当 $x > at$ 时：
    \begin{equation}
        u(x, t) = \frac{1}{2}[\varphi(x-at) + \varphi(x+at)]
    \end{equation}
    
    当 $0 < x < at$ 时：
    \begin{equation}
        u(x, t) = \frac{1}{2}[-\varphi(at-x) + \varphi(x+at)]
    \end{equation}
\end{jie}
\end{lizi}

\subsection{反射波}

固定端的边界条件导致波的\textbf{反射}。入射波到达端点后，反射为一个相位相反的波。

\begin{wulibj}[波的反射]
\begin{itemize}
    \item 固定端（$u = 0$）：反射波与入射波相位相反
    \item 自由端（$\frac{\partial u}{\partial x} = 0$）：反射波与入射波相位相同
\end{itemize}
\end{wulibj}

% ========== 第四节 ==========
\section{有界弦的振动}

对于有限长弦，需要考虑端点的多次反射。这类问题通常用分离变量法\index{分离变量法}求解更方便（下一章介绍）。

但也可以用行波叠加的思想理解：有界弦的振动可以看作初始波在两端点间来回反射叠加的结果。

% ========== 本章小结 ==========
\section{本章小结}

\subsection{核心内容}

\begin{enumerate}
    \item \textbf{达朗贝尔公式}：$u(x, t) = \frac{1}{2}[\varphi(x-at) + \varphi(x+at)] + \frac{1}{2a}\int_{x-at}^{x+at}\psi(s)\,\mathrm{d}s$
    \item \textbf{通解形式}：$u = f(x-at) + g(x+at)$（左、右行波叠加）
    \item \textbf{特征线}：$x \mp at = C$
\end{enumerate}

\subsection{关键概念}

\begin{itemize}
    \item 依赖区间：$(x, t)$ 处的解依赖于初始条件在 $[x-at, x+at]$ 上的值
    \item 影响区域：初始扰动影响的时空区域
    \item 有限传播速度：波以速度 $a$ 传播
\end{itemize}

% ========== 习题 ==========
\section{习题}

\textbf{基础题}

\begin{enumerate}
    \item 用达朗贝尔公式求解：
    \begin{equation}
        \begin{cases}
            u_{tt} = 4u_{xx} \\
            u(x, 0) = \sin x, \quad u_t(x, 0) = 0
        \end{cases}
    \end{equation}
    
    \item 验证 $u(x, t) = f(x-at) + g(x+at)$ 满足波动方程。
    
    \item 设初始位移 $\varphi(x)$ 只在 $[a, b]$ 内非零，求 $u(x, t)$ 非零的时空区域。
\end{enumerate}

\textbf{综合题}

\begin{enumerate}[resume]
    \item 求解半无界问题：
    \begin{equation}
        \begin{cases}
            u_{tt} = a^2 u_{xx}, & x > 0, \; t > 0 \\
            u(x, 0) = \mathrm{e}^{-x}, \quad u_t(x, 0) = 0 \\
            u_x(0, t) = 0 \quad \text{（自由端）}
        \end{cases}
    \end{equation}
    
    \item 设弦在 $t = 0$ 时静止，初始位移为三角形脉冲：
    \begin{equation}
        \varphi(x) = \begin{cases}
            x, & 0 \leq x \leq 1 \\
            2-x, & 1 < x \leq 2 \\
            0, & \text{其他}
        \end{cases}
    \end{equation}
    求 $u(x, t)$ 并画出 $t = 0, \frac{1}{a}, \frac{2}{a}, \frac{3}{a}$ 时的波形。
\end{enumerate}

\textbf{挑战题}

\begin{enumerate}[resume]
    \item 推导三维波动方程初值问题的解（泊松公式或基尔霍夫公式）：
    \begin{equation}
        u(\boldsymbol{x}, t) = \frac{1}{4\pi a}\frac{\partial}{\partial t}\left[\frac{1}{t}\int_{|\boldsymbol{y}|=at} \varphi(\boldsymbol{x}+\boldsymbol{y})\,\mathrm{d}S\right] + \frac{1}{4\pi at}\int_{|\boldsymbol{y}|=at} \psi(\boldsymbol{x}+\boldsymbol{y})\,\mathrm{d}S
    \end{equation}
    
    \item 比较一维、二维、三维波动方程解的性质差异（如惠更斯原理的强弱形式）。
\end{enumerate}
